\chapter{Wprowadzenie}

Dynamiczny rozwój technologii generatywnej sztucznej inteligencji doprowadził do sytuacji, w której tworzenie fałszywych treści multimedialnych stało się dostępne dla szerokiego grona użytkowników bez specjalistycznej wiedzy technicznej~\cite{Khan2025ComprehensiveSurvey,Li2025AIGCSurvey}.
Narzędzia takie jak generatory wideo oparte na modelach dyfuzyjnych, systemy klonowania głosu czy generatory obrazów umożliwiają produkcję materiałów audiowizualnych o wysokim stopniu realizmu, trudnych do odróżnienia od autentycznych nagrań.
Zjawisko to stanowi poważne zagrożenie dla bezpieczeństwa informacyjnego, integralności debaty publicznej oraz zaufania do mediów~\cite{Europol2023Deepfakes}.

Dezinformacja multimedialna -- rozumiana jako celowe rozpowszechnianie fałszywych lub zmanipulowanych treści wizualnych i dźwiękowych -- nie jest zjawiskiem nowym.
Jednak postęp w dziedzinie uczenia głębokiego, w szczególności modeli GAN (ang.~\textit{Generative Adversarial Network}) oraz dyfuzji latentnej, znacząco obniżył próg wejścia dla twórców fałszywych materiałów i jednocześnie podniósł ich jakość do poziomu, który utrudnia wykrycie przy użyciu tradycyjnych metod analizy.
Zjawisko powszechnie znane jako \textit{deepfake} ewoluowało od prostych podmian twarzy do kompleksowych syntez wideo, w których głos, mimika i kontekst narracyjny są generowane w całości przez modele sztucznej inteligencji.

W odpowiedzi na te zagrożenia środowisko naukowe i przemysłowe intensywnie pracuje nad narzędziami detekcji fałszywych treści multimedialnych.
Metody te obejmują zarówno klasyczne techniki analizy sygnału -- takie jak transformata Fouriera czy analiza przepływu optycznego -- jak i nowoczesne podejścia oparte na sieciach neuronowych, w tym konwolucyjnych sieciach CNN (ang.~\textit{Convolutional Neural Network}) i transformerach wizyjnych ViT (ang.~\textit{Vision Transformer}).
Różnorodność stosowanych podejść, brak jednolitych standardów ewaluacji oraz szybkie tempo ewolucji technik generatywnych sprawiają, że ocena skuteczności dostępnych narzędzi pozostaje otwartym i istotnym problemem badawczym~\cite{Wang2024ReliabilityPerspective}.

\section*{Cel i zakres pracy}

Celem niniejszej pracy magisterskiej jest ocena, czy wielosygnałowe podejście do detekcji wideo generowanego przez AI może pełnić rolę konserwatywnego narzędzia wspomagającego weryfikację treści w scenariuszach newsroomowych i moderacyjnych.
Praca obejmuje ukierunkowany przegląd literatury naukowej i materiałów kontekstowych, praktyczną implementację detektora oraz jego ewaluację na pilotażowym benchmarku zawierającym materiały wideo wygenerowane przez modele AI, ich warianty po kompresji i przycięciu oraz negatywy broadcastowe.
Zakres został świadomie zawężony do materiałów \textbf{wideo}; praca nie obejmuje obrazów statycznych, audio ani klasycznych face-swap deepfake'ów w nagraniach rzeczywistych.

W ramach pracy zrealizowano następujące cele szczegółowe:
\begin{enumerate}
    \item Analiza zagrożeń związanych z manipulacją treściami multimedialnymi -- taksonomia technik, skala zjawiska oraz konsekwencje społeczne.
    \item Krytyczny przegląd dostępnych metod i narzędzi detekcji wideo AI -- rozwiązań komercyjnych, open-source oraz podejść akademickich.
    \item Implementacja i ewaluacja systemu detekcji wideo generowanego przez AI, opartego na fuzji artefaktów wizualnych, przepływu optycznego, cech częstotliwościowych, sygnałów stylometrycznych oraz metadanych proweniencji.
    \item Ocena kompromisu między recall a FPR na benchmarku pilotażowym, ze szczególnym naciskiem na ograniczenie fałszywych alarmów w materiałach broadcastowych.
\end{enumerate}

\section*{Hipoteza badawcza i pytania badawcze}

Hipoteza badawcza pracy brzmi następująco: \textit{wielosygnałowy detektor wideo generowanego przez AI, łączący sygnały obrazowe, temporalne i proweniencyjne, może na pilotażowym benchmarku osiągnąć wskaźnik fałszywych alarmów nie wyższy niż $1/7 \approx 0{,}143$ dla klasy negatywnej broadcast przy jednoczesnym recall co najmniej 0,55 dla bazowego splitu \texttt{ai\_baseline}.}

Weryfikacji tej hipotezy służą trzy pytania badawcze:
\begin{enumerate}
    \item Jakie ograniczenia mają publicznie dostępne narzędzia detekcji wideo AI w warunkach zbliżonych do rzeczywistej pracy redakcyjnej?
    \item Czy ręcznie projektowana fuzja wielu sygnałów pozwala obniżyć FPR bez całkowitej utraty recall na materiałach AI?
    \item Które transformacje post-processingu (kompresja, przycięcie) najsilniej degradują skuteczność detektora i które moduły są na nie najbardziej wrażliwe?
\end{enumerate}

\noindent\textbf{Wkład własny autora.}
Indywidualny wkład autora obejmuje zaprojektowanie i implementację własnego
wielosygnałowego systemu detekcji materiałów wideo generowanych przez AI,
przygotowanie środowiska badawczego i zbioru ewaluacyjnego, opracowanie
heurystyk fuzji sygnałów oraz przeprowadzenie ilościowej i jakościowej analizy wyników.
Wkład własny obejmuje również implementację interfejsu użytkownika,
mechanizmów raportowania oraz eksperymenty porównawcze z wybranymi narzędziami komercyjnymi.

\section*{Struktura pracy}

Praca składa się z sześciu rozdziałów.
Rozdział~1 stanowi wprowadzenie do problematyki dezinformacji multimedialnej i zarysowuje cel oraz zakres pracy.
Rozdział~2 zawiera przegląd literatury naukowej dotyczącej technik generowania i detekcji fałszywych treści multimedialnych.
Rozdział~3 opisuje taksonomię zagrożeń -- techniki manipulacji obrazem, dźwiękiem i wideo, ze szczególnym uwzględnieniem metod opartych na generatywnej AI.
Rozdział~4 prezentuje przegląd narzędzi i metod detekcji, zarówno komercyjnych, jak i akademickich.
Rozdział~5 opisuje implementację i ewaluację opracowanego systemu detekcji wraz z analizą uzyskanych wyników.
Rozdział~6 zawiera wnioski końcowe oraz kierunki dalszych badań.

Uzyskane wyniki oraz zaproponowany benchmark stanowią punkt wyjścia do dalszych prac badawczych i inżynierskich nad detekcją materiałów wideo generowanych przez AI.
